% فصل سوم: تبدیلات شدت و پالایش مکانی
% این فایل با \input در main.tex یا chapter03_main.tex فراخوانی می شود.
% تمام 57 شکل فصل در محل بحث خود آمده اند؛ شکل 3-57 دو بخش دارد.

\chapter{تبدیلات شدت و پالایش مکانی}
\label{chap:intensity-spatial-filtering}

\providecommand{\BookFigure}[4][0.94]{%
\begin{figure}[H]
  \begin{latin}
  \centering
  \includegraphics[width=#1\textwidth,height=0.73\textheight,keepaspectratio]{#2}
  \end{latin}
  \caption{#3}
  \label{#4}
\end{figure}}

\section{زمینه و منطق پردازش در حوزه مکانی}

در حوزه مکانی، عملیات مستقیما روی پیکسل های تصویر انجام می شود. اگر \(f\)
تصویر ورودی، \(g\) تصویر خروجی و \(T\) عملگری باشد که روی همسایگی نقطه
\((x,y)\) اثر می کند، مدل کلی فصل چنین است:
\begin{equation}
  g(x,y)=T\!\left[f(x,y)\right].
  \tag{3-1}\label{eq:3-1}
\end{equation}
رابطه \eqref{eq:3-1} هم پردازش تک پیکسلی و هم پالایش همسایگی را پوشش
می دهد. تفاوت این دو فقط در اندازه ناحیه ای است که \(T\) می بیند.

\BookFigure[0.72]{chapter03_figures/figure_3_01.png}
{قرار گرفتن یک همسایگی \(3\times3\) روی پیکسل هدف. مرکز هسته از پیکسل به
پیکسل حرکت می کند و در هر مکان یک مقدار خروجی می سازد.}
{fig:3-1}

اگر همسایگی به یک پیکسل کاهش یابد، شدت ورودی \(r\) مستقیما به شدت خروجی
\(s\) نگاشت می شود:
\begin{equation}
  s=T(r).
  \tag{3-2}\label{eq:3-2}
\end{equation}
در این حالت، هیچ اطلاعات مکانی از پیکسل های مجاور استفاده نمی شود؛ بنابراین
می توان \(T\) را به شکل یک جدول جست وجو نیز پیاده سازی کرد.

\BookFigure[0.83]{chapter03_figures/figure_3_02.png}
{دو نگاشت پایه شدت: کشش کنتراست و آستانه گذاری. شیب تابع نشان می دهد یک
بازه شدت فشرده یا گسترده می شود.}
{fig:3-2}

\section{توابع پایه تبدیل شدت}

\subsection{خانواده تبدیلات و تصویر منفی}

شکل \ref{fig:3-3} تصویری فشرده از رفتار نگاشت های مهم این بخش است. منحنی
بالای خط همانی، مقدار خروجی را از ورودی روشن تر و منحنی پایین آن، خروجی را
تیره تر می کند.

\BookFigure[0.74]{chapter03_figures/figure_3_03.png}
{شکل کلی تبدیلات منفی، لگاریتمی، توانی و همانی در بازه شدت
\([0,L-1]\).}
{fig:3-3}

تبدیل منفی برای تصویری با \(L\) سطح شدت برابر است با
\begin{equation}
  s=L-1-r.
  \tag{3-3}\label{eq:3-3}
\end{equation}
این رابطه ترتیب روشنایی را وارونه می کند. جزئیات روشن محصور در زمینه تیره
گاهی در نسخه منفی خواناتر دیده می شوند.

\BookFigure[0.82]{chapter03_figures/figure_3_04.png}
{ماموگرام اصلی و منفی آن. اطلاعات مکانی یکسان است، اما ترتیب شدت ها وارونه
شده است.}
{fig:3-4}

\subsection{تبدیل لگاریتمی}

فرم عمومی تبدیل لگاریتمی چنین است:
\begin{equation}
  s=c\log(1+r), \qquad r\ge 0.
  \tag{3-4}\label{eq:3-4}
\end{equation}
افزودن ۱ از محاسبه لگاریتم صفر جلوگیری می کند. این تبدیل اختلاف بین شدت های
کوچک را گسترش می دهد و مقادیر بزرگ را فشرده می کند؛ از این رو برای نمایش
طیف فوریه مناسب است.

\BookFigure[0.84]{chapter03_figures/figure_3_05.png}
{نمایش یک طیف فوریه پیش و پس از تبدیل لگاریتمی. مولفه های کم دامنه پس از
فشرده شدن دامنه دینامیکی آشکار می شوند.}
{fig:3-5}

\subsection{تبدیل توانی و تصحیح گاما}

تبدیل توانی یا گاما با رابطه زیر تعریف می شود:
\begin{equation}
  s=c r^{\gamma}, \qquad c>0,\quad \gamma>0.
  \tag{3-5}\label{eq:3-5}
\end{equation}
برای ورودی نرمال شده، \(\gamma<1\) ناحیه های تیره را باز می کند و
\(\gamma>1\) آنها را فشرده می سازد. مقدار \(\gamma=1\) نگاشت همانی است.

\BookFigure[0.72]{chapter03_figures/figure_3_06.png}
{اثر مقادیر مختلف \(\gamma\) بر منحنی تبدیل توانی. فاصله منحنی از خط همانی،
میزان تغییر روشنایی را نشان می دهد.}
{fig:3-6}

\BookFigure[0.86]{chapter03_figures/figure_3_07.png}
{ایده تصحیح گاما: کدگذاری معکوس پاسخ غیرخطی نمایشگر سبب می شود خروجی دیده
شده به رمپ اصلی نزدیک شود.}
{fig:3-7}

\BookFigure[0.70]{chapter03_figures/figure_3_08.png}
{تصویر پزشکی و نتایج گاماهای کمتر از یک. با کاهش گاما، جزئیات تیره روشن تر
می شوند، اما گامای بسیار کوچک ظاهر تصویر را می شوید.}
{fig:3-8}

\BookFigure[0.76]{chapter03_figures/figure_3_09.png}
{تصویر هوایی و نتایج گاماهای بزرگ تر از یک. انتخاب گاما باید با دامنه شدت
تصویر و هدف نمایش هماهنگ باشد.}
{fig:3-9}

\subsection{توابع چندپاره، برش شدت و صفحه های بیت}

در تبدیل چندپاره، هر بازه ورودی شیب مستقلی دارد. شیب بزرگ تر از یک کنتراست
آن بازه را افزایش می دهد و شیب کوچک تر از یک آن را فشرده می کند. آستانه گذاری
حد نهایی همین ایده است.

\BookFigure[0.78]{chapter03_figures/figure_3_10.png}
{تابع خطی چندپاره برای کشش کنتراست، تصویر کم کنتراست، نتیجه کشش و نتیجه
آستانه گذاری.}
{fig:3-10}

\BookFigure[0.80]{chapter03_figures/figure_3_11.png}
{دو گونه برش شدت: در حالت نخست شدت های بیرون بازه سرکوب می شوند؛ در حالت
دوم زمینه دست نخورده می ماند.}
{fig:3-11}

\BookFigure[0.90]{chapter03_figures/figure_3_12.png}
{برش شدت در آنژیوگرام. انتخاب نوع تابع تعیین می کند زمینه حذف شود یا همراه
رگ ها باقی بماند.}
{fig:3-12}

در تصویر ۸ بیتی، مقدار هر پیکسل را می توان به صورت
\(r=\sum_{j=0}^{7}b_j2^j\) نوشت. صفحه های پرارزش ساختار کلی و صفحه های
کم ارزش جزئیات ظریف و نویز را بیشتر حمل می کنند.

\BookFigure[0.78]{chapter03_figures/figure_3_13.png}
{تجزیه یک تصویر ۸ بیتی به هشت صفحه تک بیتی، از بیت پرارزش تا بیت کم ارزش.}
{fig:3-13}

\BookFigure[0.92]{chapter03_figures/figure_3_14.png}
{تصویر اصلی و هشت صفحه بیت آن. سهم ساختاری بیت های بالا و سهم جزئیات ریز
بیت های پایین به صورت دیداری قابل مقایسه است.}
{fig:3-14}

\BookFigure[0.92]{chapter03_figures/figure_3_15.png}
{بازسازی تصویر فقط با صفحه های پرارزش. افزودن هر صفحه بیت، بخشی از جزئیات
شدت را بازمی گرداند.}
{fig:3-15}

\begin{latin}
\subsubsection*{Executable experiment: pointwise intensity transformations}
The script applies the negative, logarithmic, gamma, and threshold transforms to
the same deterministic sample. Keeping the input fixed makes the visual effect of
each mapping directly comparable.
\lstinputlisting[caption={Pointwise intensity transformations},label={code:intensity-transform}]
{chapter03_code/intensity_transform_demo.py}
\begin{figure}[H]
  \centering
  \includegraphics[width=0.96\textwidth]{chapter03_code_outputs/intensity_transform_result.png}
  \caption*{Executed output of the pointwise intensity transformation script}
\end{figure}
\end{latin}

\section{پردازش هیستوگرام}

\subsection{تعریف هیستوگرام و خواندن شکل آن}

اگر \(r_k\) شدت \(k\)-ام و \(n_k\) تعداد پیکسل های دارای آن شدت باشد،
هیستوگرام نرمال نشده و نرمال شده به ترتیب عبارت اند از
\begin{align}
  h(r_k)&=n_k,\qquad k=0,1,\ldots,L-1,
  \tag{3-6}\label{eq:3-6}\\
  p_r(r_k)&=\frac{n_k}{MN},\qquad k=0,1,\ldots,L-1.
  \tag{3-7}\label{eq:3-7}
\end{align}
مجموع \(p_r(r_k)\) برابر یک است؛ پس می توان آن را تقریب گسسته تابع چگالی
احتمال شدت دانست.

\BookFigure[0.95]{chapter03_figures/figure_3_16.png}
{چهار الگوی اصلی هیستوگرام: تصویر تیره، روشن، کم کنتراست و پرکنتراست. محل
و پهنای جرم هیستوگرام، روشنایی و کنتراست را خلاصه می کند.}
{fig:3-16}

برای همسان سازی پیوسته، نگاشت باید افزایشی باشد تا ترتیب شدت ها حفظ شود:
\begin{align}
  s&=T(r),\qquad 0\le r\le L-1,
  \tag{3-8}\label{eq:3-8}\\
  r&=T^{-1}(s),\qquad 0\le s\le L-1.
  \tag{3-9}\label{eq:3-9}
\end{align}

\BookFigure[0.80]{chapter03_figures/figure_3_17.png}
{تفاوت تابع افزایشی و تابع اکیدا افزایشی. فقط حالت دوم نگاشت یک به یک و
معکوس یکتا می دهد.}
{fig:3-17}

قانون تغییر متغیر برای چگالی ها چنین است:
\begin{equation}
  p_s(s)=p_r(r)\left|\frac{dr}{ds}\right|.
  \tag{3-10}\label{eq:3-10}
\end{equation}
انتخاب \(T\) به صورت تابع توزیع تجمعی، هیستوگرام خروجی را در حالت پیوسته
یکنواخت می کند:
\begin{equation}
  s=T(r)=(L-1)\int_0^r p_r(w)\,dw.
  \tag{3-11}\label{eq:3-11}
\end{equation}
مشتق این تبدیل برابر است با
\begin{equation}
  \frac{ds}{dr}=\frac{d}{dr}\left[(L-1)\int_0^r p_r(w)\,dw\right]
  =(L-1)p_r(r).
  \tag{3-12}\label{eq:3-12}
\end{equation}
با جای گذاری در رابطه \eqref{eq:3-10} به نتیجه زیر می رسیم:
\begin{equation}
  p_s(s)=\frac{1}{L-1},\qquad 0\le s\le L-1.
  \tag{3-13}\label{eq:3-13}
\end{equation}

\BookFigure[0.86]{chapter03_figures/figure_3_18.png}
{در مدل پیوسته، اعمال تابع توزیع تجمعی به هر چگالی ورودی، چگالی یکنواخت
خروجی می دهد.}
{fig:3-18}

در تصویر دیجیتال، احتمال هر سطح همان نسبت شمارش پیکسل ها است:
\begin{equation}
  p_r(r_k)=\frac{n_k}{MN}.
  \tag{3-14}\label{eq:3-14}
\end{equation}
نسخه گسسته همسان سازی هیستوگرام از جمع تجمعی استفاده می کند:
\begin{equation}
  s_k=T(r_k)=(L-1)\sum_{j=0}^{k}p_r(r_j),
  \qquad k=0,1,\ldots,L-1.
  \tag{3-15}\label{eq:3-15}
\end{equation}
چون \(s_k\) ممکن است کسری باشد، باید به نزدیک ترین سطح مجاز گرد شود. برخلاف
حالت پیوسته، یکنواخت شدن کامل تضمین نمی شود؛ چند سطح ورودی ممکن است به یک
سطح خروجی نگاشت شوند.

\begin{table}[H]
\centering
\caption{توزیع شدت یک تصویر ۳ بیتی با اندازه \(64\times64\)}
\label{tab:3-1}
\begin{tabular}{ccc}
\toprule
\(r_k\) & \(n_k\) & \(p_r(r_k)=n_k/(MN)\)\\
\midrule
0 & 790  & 0.19\\
1 & 1023 & 0.25\\
2 & 850  & 0.21\\
3 & 656  & 0.16\\
4 & 329  & 0.08\\
5 & 245  & 0.06\\
6 & 122  & 0.03\\
7 & 81   & 0.02\\
\bottomrule
\end{tabular}
\end{table}

\BookFigure[0.90]{chapter03_figures/figure_3_19.png}
{مثال گسسته همسان سازی: هیستوگرام ورودی، تابع تجمعی پله ای و هیستوگرام
خروجی. ادغام چند سطح علت تخت نبودن خروجی است.}
{fig:3-19}

معکوس گسسته به صورت
\begin{equation}
  r_k=T^{-1}(s_k)
  \tag{3-16}\label{eq:3-16}
\end{equation}
نوشته می شود. این معکوس فقط وقتی یکتا است که نگاشت اکیدا افزایشی باشد.

\BookFigure[0.78]{chapter03_figures/figure_3_20.png}
{اثر همسان سازی سراسری بر چهار تصویر دارای محتوای یکسان و کنتراست متفاوت.
ستون سوم، هیستوگرام خروجی هر مورد را نشان می دهد.}
{fig:3-20}

\BookFigure[0.82]{chapter03_figures/figure_3_21.png}
{توابع تبدیل چهار تصویر شکل پیشین. تابع نزدیک به خط همانی یعنی ورودی از قبل
دامنه شدت مناسبی داشته است.}
{fig:3-21}

\subsection{تطبیق هیستوگرام}

اگر هدف، هیستوگرام یکنواخت نباشد و چگالی مطلوب \(p_z(z)\) مشخص شود، ابتدا
ورودی را با
\begin{equation}
  s=T(r)=(L-1)\int_0^r p_r(w)\,dw
  \tag{3-17}\label{eq:3-17}
\end{equation}
به متغیر میانی \(s\) می بریم. سپس برای خروجی مطلوب تعریف می کنیم
\begin{equation}
  G(z)=(L-1)\int_0^z p_z(v)\,dv=s,
  \tag{3-18}\label{eq:3-18}
\end{equation}
و در نتیجه نگاشت نهایی برابر است با
\begin{equation}
  z=G^{-1}(s)=G^{-1}[T(r)].
  \tag{3-19}\label{eq:3-19}
\end{equation}

در حالت گسسته، سه گام متناظر چنین اند:
\begin{align}
  s_k&=T(r_k)=(L-1)\sum_{j=0}^{k}p_r(r_j),
  \tag{3-20}\label{eq:3-20}\\
  G(z_q)&=(L-1)\sum_{i=0}^{q}p_z(z_i),
  \tag{3-21}\label{eq:3-21}\\
  G(z_q)&=s_k,
  \tag{3-22}\label{eq:3-22}\\
  z_q&=G^{-1}(s_k).
  \tag{3-23}\label{eq:3-23}
\end{align}
در پیاده سازی، برای هر \(s_k\) نزدیک ترین مقدار \(G(z_q)\) انتخاب می شود؛
در تساوی فاصله، کوچک ترین \(z_q\) انتخاب ثابتی و قابل بازتولید می دهد.

\BookFigure[0.82]{chapter03_figures/figure_3_22.png}
{مثال ۳ بیتی تطبیق هیستوگرام: ورودی، هیستوگرام مطلوب، تابع تجمعی مطلوب و
هیستوگرام واقعی خروجی.}
{fig:3-22}

\begin{table}[H]
\centering
\caption{هیستوگرام مطلوب و هیستوگرام واقعی مثال ۳ بیتی}
\label{tab:3-2}
\begin{tabular}{ccc}
\toprule
\(z_q\) & \(p_z(z_q)\) مطلوب & \(p_z(z_q)\) واقعی\\
\midrule
0 & 0.00 & 0.00\\
1 & 0.00 & 0.00\\
2 & 0.00 & 0.00\\
3 & 0.15 & 0.19\\
4 & 0.20 & 0.25\\
5 & 0.30 & 0.21\\
6 & 0.20 & 0.24\\
7 & 0.15 & 0.11\\
\bottomrule
\end{tabular}
\end{table}

\begin{table}[H]
\centering
\caption{مقادیر گرد شده تابع \(G(z_q)\)}
\label{tab:3-3}
\begin{tabular}{cc@{\qquad}cc}
\toprule
\(z_q\) & \(G(z_q)\) & \(z_q\) & \(G(z_q)\)\\
\midrule
0 & 0 & 4 & 2\\
1 & 0 & 5 & 5\\
2 & 0 & 6 & 6\\
3 & 1 & 7 & 7\\
\bottomrule
\end{tabular}
\end{table}

\begin{table}[H]
\centering
\caption{نگاشت مقادیر همسان شده به مقادیر خروجی مطلوب}
\label{tab:3-4}
\begin{tabular}{cc}
\toprule
\(s_k\) & \(z_q\)\\
\midrule
1 & 3\\
3 & 4\\
5 & 5\\
6 & 6\\
7 & 7\\
\bottomrule
\end{tabular}
\end{table}

\BookFigure[0.86]{chapter03_figures/figure_3_23.png}
{تصویر فوبوس و هیستوگرام متمرکز آن در شدت های پایین. تمرکز جرم نزدیک صفر
هشداری است که همسان سازی ممکن است بیش از حد روشن کند.}
{fig:3-23}

\BookFigure[0.82]{chapter03_figures/figure_3_24.png}
{تابع همسان سازی، تصویر بیش از حد روشن شده و هیستوگرام آن. جهش تند تابع در
ابتدای بازه، علت شسته شدن تصویر است.}
{fig:3-24}

\BookFigure[0.80]{chapter03_figures/figure_3_25.png}
{تطبیق هیستوگرام برای همان تصویر. شکل هدف مناسب، جرم شدت را کنترل شده تر
جابجا می کند و از روشن شدن افراطی می کاهد.}
{fig:3-25}

\subsection{پردازش محلی و آمار همسایگی}

همسان سازی سراسری یک نگاشت مشترک برای همه پیکسل ها دارد، اما در پردازش
محلی، هیستوگرام در پنجره ای پیرامون هر پیکسل محاسبه می شود.

\BookFigure[0.94]{chapter03_figures/figure_3_26.png}
{مقایسه همسان سازی سراسری و محلی. روش محلی جزئیات کم کنتراست را آشکار می کند،
اما ممکن است نویز را نیز تقویت کند.}
{fig:3-26}

گشتاور مرتبه \(n\) حول میانگین، میانگین و واریانس به ترتیب عبارت اند از
\begin{align}
  \mu_n&=\sum_{i=0}^{L-1}(r_i-m)^n p(r_i),
  \tag{3-24}\label{eq:3-24}\\
  m&=\sum_{i=0}^{L-1}r_i p(r_i),
  \tag{3-25}\label{eq:3-25}\\
  \sigma^2=\mu_2&=\sum_{i=0}^{L-1}(r_i-m)^2p(r_i).
  \tag{3-26}\label{eq:3-26}
\end{align}
میانگین شاخص روشنایی متوسط و واریانس شاخص پراکندگی شدت یا کنتراست است.

برای همسایگی \(S_{xy}\) مرکز شده روی \((x,y)\)، آمار محلی چنین تعریف می شود:
\begin{align}
  m_{S_{xy}}&=\sum_{i=0}^{L-1}r_i p_{S_{xy}}(r_i),
  \tag{3-27}\label{eq:3-27}\\
  \sigma^2_{S_{xy}}&=\sum_{i=0}^{L-1}(r_i-m_{S_{xy}})^2p_{S_{xy}}(r_i).
  \tag{3-28}\label{eq:3-28}
\end{align}

یک قانون انتخابی برای تقویت ناحیه های تیره و کم کنتراست به صورت زیر است:
\begin{equation}
g(x,y)=
\begin{cases}
C f(x,y), & k_0m_G\le m_{S_{xy}}\le k_1m_G
\;\land\; k_2\sigma_G\le\sigma_{S_{xy}}\le k_3\sigma_G,\\
f(x,y), & \text{در غیر این صورت.}
\end{cases}
\tag{3-29}\label{eq:3-29}
\end{equation}
مزیت رابطه \eqref{eq:3-29} این است که فقط پیکسل های واجد دو شرط روشنایی و
کنتراست تغییر می کنند؛ بنابراین زمینه و ناحیه های از قبل مناسب حفظ می شوند.

\BookFigure[0.82]{chapter03_figures/figure_3_27.png}
{تقویت انتخابی بر اساس میانگین و انحراف معیار محلی. ساختارهای پنهان آشکار
می شوند، در حالی که زمینه کمتر از همسان سازی محلی تغییر می کند.}
{fig:3-27}

\begin{latin}
\subsubsection*{Executable experiment: global and local histogram processing}
The script first compresses the input contrast, then applies global histogram
equalization and a transparent tile-based local equalization. Each output is shown
with its measured histogram.
\lstinputlisting[caption={Global and local histogram processing},label={code:histogram-processing}]
{chapter03_code/histogram_processing_demo.py}
\begin{figure}[H]
  \centering
  \includegraphics[width=0.98\textwidth]{chapter03_code_outputs/histogram_processing_result.png}
  \caption*{Executed output of the histogram processing script}
\end{figure}
\end{latin}

\section{اصول پالایش مکانی}

\subsection{جمع حاصل ضرب و حرکت هسته}

برای یک هسته \(3\times3\)، پاسخ در نقطه \((x,y)\) جمع نه حاصل ضرب است:
\begin{align}
g(x,y)={}&w(-1,-1)f(x-1,y-1)+w(-1,0)f(x-1,y)+\cdots\notag\\
&+w(0,0)f(x,y)+\cdots+w(1,1)f(x+1,y+1).
\tag{3-30}\label{eq:3-30}
\end{align}
نسخه عمومی برای هسته \(m\times n\)، با \(m=2a+1\) و \(n=2b+1\)، چنین است:
\begin{equation}
  g(x,y)=\sum_{s=-a}^{a}\sum_{t=-b}^{b}w(s,t)f(x+s,y+t).
  \tag{3-31}\label{eq:3-31}
\end{equation}

\BookFigure[0.72]{chapter03_figures/figure_3_28.png}
{مکان هندسی هسته و پیکسل های زیر آن. ضریب مرکزی \(w(0,0)\) با پیکسل
\(f(x,y)\) هم مکان است.}
{fig:3-28}

در یک بعد، همبستگی به صورت
\begin{equation}
  g(x)=\sum_{s=-a}^{a}w(s)f(x+s)
  \tag{3-32}\label{eq:3-32}
\end{equation}
نوشته می شود. در هم پیچش، هسته پیش از حرکت ۱۸۰ درجه می چرخد؛ بنابراین پاسخ
ضربه، جهت هسته را در همبستگی حفظ و در هم پیچش وارونه می کند.

\BookFigure[0.82]{chapter03_figures/figure_3_29.png}
{مقایسه مرحله به مرحله همبستگی و هم پیچش یک بعدی با ورودی ضربه. تفاوت اصلی،
چرخش هسته در هم پیچش است.}
{fig:3-29}

ضربه دوبعدی در نقطه \((x_0,y_0)\) چنین تعریف می شود:
\begin{equation}
\delta(x-x_0,y-y_0)=
\begin{cases}
1,& x=x_0,\ y=y_0,\\
0,& \text{در غیر این صورت.}
\end{cases}
\tag{3-33}\label{eq:3-33}
\end{equation}

\BookFigure[0.78]{chapter03_figures/figure_3_30.png}
{همبستگی و هم پیچش دوبعدی یک هسته نامتقارن با ضربه. خروجی همبستگی نسخه
چرخیده و خروجی هم پیچش نسخه هم جهت هسته را نشان می دهد.}
{fig:3-30}

نمادگذاری فشرده همبستگی و هم پیچش به ترتیب چنین است:
\begin{align}
  (w\star f)(x,y)&=\sum_{s=-a}^{a}\sum_{t=-b}^{b}w(s,t)f(x+s,y+t),
  \tag{3-34}\label{eq:3-34}\\
  (w*f)(x,y)&=\sum_{s=-a}^{a}\sum_{t=-b}^{b}w(s,t)f(x-s,y-t).
  \tag{3-35}\label{eq:3-35}
\end{align}
اگر تصویر \(M\times N\) و هسته \(m\times n\) باشد، اندازه خروجی کامل برابر است با
\begin{align}
  S_v&=m+M-1,
  \tag{3-36}\label{eq:3-36}\\
  S_h&=n+N-1.
  \tag{3-37}\label{eq:3-37}
\end{align}
خروجی هم اندازه با تصویر، بخش مرکزی این خروجی کامل است و به قرارداد پدگذاری
وابسته است.

\begin{latin}
\subsubsection*{Executable experiment: correlation versus convolution}
An asymmetric kernel is essential in this experiment. A symmetric kernel would
hide the 180-degree rotation and make the two outputs appear identical.
\lstinputlisting[caption={Correlation and convolution of a unit impulse},label={code:corr-conv}]
{chapter03_code/correlation_convolution_demo.py}
\begin{figure}[H]
  \centering
  \includegraphics[width=0.98\textwidth]{chapter03_code_outputs/correlation_convolution_result.png}
  \caption*{Executed output showing the orientation difference between correlation and convolution}
\end{figure}
\end{latin}

\subsection{خواص، هسته مرکب و جدایی پذیری}

اگر پالایش در \(Q\) مرحله با هسته های \(w_1,\ldots,w_Q\) انجام شود، هسته
معادل یک مرحله ای برابر است با
\begin{equation}
  w=w_1*w_2*\cdots*w_Q.
  \tag{3-38}\label{eq:3-38}
\end{equation}
برای هسته های هم اندازه \(m\times n\)، اندازه هسته مرکب چنین است:
\begin{align}
  W_v&=Q(m-1)+m,
  \tag{3-39}\label{eq:3-39}\\
  W_h&=Q(n-1)+n.
  \tag{3-40}\label{eq:3-40}
\end{align}

\begin{table}[H]
\centering
\caption{خواص پایه هم پیچش و همبستگی}
\label{tab:3-5}
\begin{tabular}{p{0.22\textwidth}p{0.34\textwidth}p{0.34\textwidth}}
\toprule
خاصیت & هم پیچش & همبستگی\\
\midrule
جابجایی & \(f*g=g*f\) & برقرار نیست\\
شرکت پذیری & \(f*(g*h)=(f*g)*h\) & برقرار نیست\\
توزیع پذیری & \(f*(g+h)=f*g+f*h\) & \(f\star(g+h)=f\star g+f\star h\)\\
\bottomrule
\end{tabular}
\end{table}

\BookFigure[0.84]{chapter03_figures/figure_3_31.png}
{هسته میانگین گیری جعبه ای و هسته گاوسی \(3\times3\)، هر دو نرمال شده با
مجموع ضرایب یک.}
{fig:3-31}

هسته جدایی پذیر حاصل ضرب خارجی دو بردار است:
\begin{align}
  w&=vw^{T},
  \tag{3-41}\label{eq:3-41}\\
  w&=vv^{T}\qquad\text{برای هسته مربعی متقارن.}
  \tag{3-42}\label{eq:3-42}
\end{align}
به کمک شرکت پذیری هم پیچش داریم
\begin{equation}
w*f=(w_1*w_2)*f=w_2*(w_1*f).
\tag{3-43}\label{eq:3-43}
\end{equation}
پس پالایش دوبعدی با دو پالایش یک بعدی جایگزین می شود. نسبت تقریبی هزینه روش
دوبعدی به روش جدایی پذیر برابر است با
\begin{equation}
  C=\frac{MNmn}{MN(m+n)}=\frac{mn}{m+n}.
  \tag{3-44}\label{eq:3-44}
\end{equation}
برای هسته های بزرگ، این کاهش محاسبه بسیار مهم است. رتبه یک بودن ماتریس هسته،
آزمون مستقیم جدایی پذیری است.

\BookFigure[0.88]{chapter03_figures/figure_3_32.png}
{رابطه حوزه فرکانس و مکان برای یک پایین گذر ایده آل. قطع ناگهانی فرکانس، دنباله
نوسانی در هسته مکانی تولید می کند.}
{fig:3-32}

\section{پالایه های مکانی هموارساز}

\subsection{هسته جعبه ای، اندازه هسته و پدگذاری}

هسته جعبه ای \(m\times n\) همه ضرایب را برابر \(1/(mn)\) می گیرد. این
نرمال سازی شدت ثابت را ثابت نگه می دارد، زیرا مجموع ضرایب یک است.

\BookFigure[0.72]{chapter03_figures/figure_3_33.png}
{افزایش اندازه هسته جعبه ای، تاری را بیشتر و جزئیات ریز را کمتر می کند. جهت
های افقی و عمودی هسته در نتیجه قابل مشاهده اند.}
{fig:3-33}

\BookFigure[0.78]{chapter03_figures/figure_3_34.png}
{فاصله هر ضریب از مرکز در هسته های مربعی. این فاصله مبنای محاسبه وزن گاوسی
است.}
{fig:3-34}

\subsection{هسته گاوسی}

هسته گاوسی دوبعدی با ثابت مقیاس \(K\) و انحراف معیار \(\sigma\) چنین است:
\begin{equation}
  w(s,t)=G(s,t)=K\exp\!\left[-\frac{s^2+t^2}{2\sigma^2}\right].
  \tag{3-45}\label{eq:3-45}
\end{equation}
با \(r=\sqrt{s^2+t^2}\)، فرم شعاعی آن می شود
\begin{equation}
  G(r)=K\exp\!\left[-\frac{r^2}{2\sigma^2}\right].
  \tag{3-46}\label{eq:3-46}
\end{equation}
وابستگی فقط به \(r\) سبب همسانگردی می شود. پس از نمونه برداری، ضرایب بر مجموع
خود تقسیم می شوند تا بهره \lr{DC} برابر یک شود.

\BookFigure[0.86]{chapter03_figures/figure_3_35.png}
{نمونه برداری از تابع گاوسی دوبعدی و تشکیل هسته نرمال شده \(3\times3\).}
{fig:3-35}

برای هسته مربعی \(m\times m\)، بیشترین فاصله در گوشه رخ می دهد:
\begin{equation}
  r_{\max}^{2}=2\left(\frac{m-1}{2}\right)^2=\frac{(m-1)^2}{2}.
  \tag{3-47}\label{eq:3-47}
\end{equation}
در عمل اندازه ای نزدیک \(6\sigma\times6\sigma\)، با گرد کردن به نزدیک ترین
عدد فرد، بخش موثر تابع گاوسی را پوشش می دهد.

\begin{table}[H]
\centering
\caption{پارامترهای حاصل ضرب و هم پیچش دو تابع گاوسی یک بعدی}
\label{tab:3-6}
\small
\setlength{\tabcolsep}{3pt}
\begin{tabular}{p{0.16\textwidth}p{0.10\textwidth}p{0.10\textwidth}p{0.23\textwidth}p{0.23\textwidth}}
\toprule
پارامتر & \(f\) & \(g\) & \(f\times g\) & \(f*g\)\\
\midrule
میانگین & \(m_f\) & \(m_g\) &
\(\dfrac{m_f\sigma_g^2+m_g\sigma_f^2}{\sigma_f^2+\sigma_g^2}\) &
\(m_f+m_g\)\\[2mm]
انحراف معیار & \(\sigma_f\) & \(\sigma_g\) &
\(\dfrac{\sigma_f\sigma_g}{\sqrt{\sigma_f^2+\sigma_g^2}}\) &
\(\sqrt{\sigma_f^2+\sigma_g^2}\)\\
\bottomrule
\end{tabular}
\end{table}

\BookFigure[0.92]{chapter03_figures/figure_3_36.png}
{مقایسه تاری گاوسی برای دو مقدار \(\sigma\). افزایش \(\sigma\) پهنای هسته موثر
و مقدار هموارسازی را افزایش می دهد.}
{fig:3-36}

\BookFigure[0.92]{chapter03_figures/figure_3_37.png}
{هسته بزرگ تر از دامنه موثر گاوسی بهبود محسوسی ایجاد نمی کند؛ تصویر تفاضل
تقریبا سیاه است.}
{fig:3-37}

\BookFigure[0.92]{chapter03_figures/figure_3_38.png}
{پروفایل لبه پس از پالایش جعبه ای و گاوسی. شکل پاسخ گذار، اثر ضرایب هسته را
مستقیم نشان می دهد.}
{fig:3-38}

\BookFigure[0.92]{chapter03_figures/figure_3_39.png}
{مقایسه پدگذاری صفر، آینه ای و تکرار مرز. پدگذاری صفر در مرز تصویر نوار تیره
مصنوعی ایجاد می کند.}
{fig:3-39}

\BookFigure[0.92]{chapter03_figures/figure_3_40.png}
{وابستگی هموارسازی به اندازه نسبی هسته و تصویر. یک هسته ثابت روی تصویر بزرگ
اثر نسبی کمتری دارد.}
{fig:3-40}

\BookFigure[0.90]{chapter03_figures/figure_3_41.png}
{هموارسازی و آستانه گذاری برای استخراج ناحیه های بزرگ و روشن از تصویر
نجومی.}
{fig:3-41}

\BookFigure[0.92]{chapter03_figures/figure_3_42.png}
{تصحیح سایه: الگوی روشنایی آهسته با پایین گذر برآورد و تصویر بر آن تقسیم
می شود.}
{fig:3-42}

\subsection{پالایش میانه}

پالایه میانه غیرخطی است: مقادیر همسایگی مرتب و مقدار میانی جایگزین پیکسل
مرکزی می شود. چون بر میانگین تکیه ندارد، برای نویز ضربه ای بسیار موثر است و
لبه ها را بهتر از میانگین گیری حفظ می کند.

\BookFigure[0.92]{chapter03_figures/figure_3_43.png}
{مقایسه گاوسی و میانه روی نویز نمک و فلفل. میانه نقاط ضربه ای را بهتر حذف
می کند و تاری کمتری به مرزها می دهد.}
{fig:3-43}

\begin{latin}
\subsubsection*{Executable experiment: smoothing filters}
The same salt-and-pepper realization is processed by box, Gaussian, and median
filters. The fixed input makes noise rejection and edge blurring directly comparable.
\lstinputlisting[caption={Box, Gaussian, and median smoothing},label={code:smoothing}]
{chapter03_code/smoothing_filter_demo.py}
\begin{figure}[H]
  \centering
  \includegraphics[width=0.99\textwidth]{chapter03_code_outputs/smoothing_filter_result.png}
  \caption*{Executed comparison of smoothing filters on one noisy image}
\end{figure}
\end{latin}

\section{پالایه های مکانی تیزکننده}

\subsection{مشتق های گسسته و لاپلاسین}

تعریف تفاضلی مشتق مرتبه اول در یک بعد عبارت است از
\begin{equation}
  \frac{\partial f}{\partial x}=f(x+1)-f(x).
  \tag{3-48}\label{eq:3-48}
\end{equation}
و مشتق مرتبه دوم به صورت
\begin{equation}
  \frac{\partial^2f}{\partial x^2}=f(x+1)+f(x-1)-2f(x)
  \tag{3-49}\label{eq:3-49}
\end{equation}
تعریف می شود. مشتق اول روی رمپ غیرصفر می ماند، اما مشتق دوم درون رمپ صفر و
در ابتدا و انتهای آن دارای علامت های مخالف است.

\BookFigure[0.84]{chapter03_figures/figure_3_44.png}
{پروفایل دارای ناحیه ثابت، رمپ و پله و پاسخ مشتق اول و دوم. عبور از صفر مشتق
دوم مکان گذار را مشخص می کند.}
{fig:3-44}

لاپلاسین جمع مشتق های مرتبه دوم در دو راستا است:
\begin{equation}
  \nabla^2 f=\frac{\partial^2f}{\partial x^2}+\frac{\partial^2f}{\partial y^2}.
  \tag{3-50}\label{eq:3-50}
\end{equation}
تقریب گسسته هر مولفه برابر است با
\begin{align}
  \frac{\partial^2f}{\partial x^2}&=f(x+1,y)+f(x-1,y)-2f(x,y),
  \tag{3-51}\label{eq:3-51}\\
  \frac{\partial^2f}{\partial y^2}&=f(x,y+1)+f(x,y-1)-2f(x,y).
  \tag{3-52}\label{eq:3-52}
\end{align}
پس لاپلاسین چهار همسایه چنین می شود:
\begin{align}
\nabla^2f(x,y)={}&f(x+1,y)+f(x-1,y)+f(x,y+1)+f(x,y-1)\notag\\
&-4f(x,y).
\tag{3-53}\label{eq:3-53}
\end{align}

\BookFigure[0.94]{chapter03_figures/figure_3_45.png}
{چهار هسته لاپلاسین. دو هسته نخست مرکز منفی و دو هسته بعدی مرکز مثبت دارند؛
این علامت در ترکیب با تصویر اصلی تعیین کننده است.}
{fig:3-45}

تیزسازی لاپلاسین با افزودن پاسخ مشتق به تصویر انجام می شود:
\begin{equation}
  g(x,y)=f(x,y)+c\nabla^2f(x,y).
  \tag{3-54}\label{eq:3-54}
\end{equation}
برای هسته با ضریب مرکزی منفی، \(c=-1\)، و برای هسته با ضریب مرکزی مثبت،
\(c=1\) است.

\BookFigure[0.70]{chapter03_figures/figure_3_46.png}
{تصویر ماه، پاسخ لاپلاسین و دو نتیجه تیزسازی. افزودن جهت های قطری جزئیات
بیشتری را تقویت می کند.}
{fig:3-46}

\BookFigure[0.55]{chapter03_figures/figure_3_47.png}
{نمایش مقیاس شده پاسخ لاپلاسین. زمینه خاکستری مقدارهای نزدیک صفر و نوارهای
روشن و تیره دو علامت مشتق دوم را نشان می دهند.}
{fig:3-47}

\subsection{ماسک ناتیز و تقویت بالاگذر}

اگر \(\bar f(x,y)\) نسخه هموار شده باشد، ماسک ناتیز برابر است با
\begin{equation}
  g_{\mathrm{mask}}(x,y)=f(x,y)-\bar f(x,y).
  \tag{3-55}\label{eq:3-55}
\end{equation}
تصویر تیز شده با ضریب بهره \(k\) از رابطه زیر به دست می آید:
\begin{equation}
  g(x,y)=f(x,y)+k g_{\mathrm{mask}}(x,y).
  \tag{3-56}\label{eq:3-56}
\end{equation}
\(k=1\) ماسک ناتیز و \(k>1\) تقویت بالاگذر است. \(k\) بزرگ هاله تیره و روشن
در اطراف لبه ها ایجاد می کند.

\BookFigure[0.84]{chapter03_figures/figure_3_48.png}
{سازوکار یک بعدی ماسک ناتیز: کم کردن نسخه هموار، مولفه لبه را جدا و افزودن
آن شیب گذار را تندتر می کند.}
{fig:3-48}

\BookFigure[0.94]{chapter03_figures/figure_3_49.png}
{تصویر متن، نسخه هموار، ماسک، نتیجه ناتیز و نتیجه تقویت بالاگذر. افزایش بهره
هم تیزی و هم هاله را بیشتر می کند.}
{fig:3-49}

\subsection{گرادیان و عملگرهای رابرتس و سوبل}

گرادیان تصویر بردار مشتق های مرتبه اول است:
\begin{equation}
  \nabla f=\operatorname{grad}(f)=
  \begin{bmatrix}g_x\\g_y\end{bmatrix}=
  \begin{bmatrix}\partial f/\partial x\\\partial f/\partial y\end{bmatrix}.
  \tag{3-57}\label{eq:3-57}
\end{equation}
اندازه گرادیان برابر است با
\begin{equation}
  M(x,y)=\|\nabla f\|=\sqrt{g_x^2+g_y^2}.
  \tag{3-58}\label{eq:3-58}
\end{equation}
تقریب سریع تر ولی کمتر همسانگرد آن چنین است:
\begin{equation}
  M(x,y)\approx |g_x|+|g_y|.
  \tag{3-59}\label{eq:3-59}
\end{equation}

\BookFigure[0.80]{chapter03_figures/figure_3_50.png}
{شماره گذاری همسایگی \(3\times3\)، دو هسته رابرتس و دو هسته سوبل. مجموع
ضرایب هر هسته مشتق صفر است.}
{fig:3-50}

تفاضل های ضربدری رابرتس عبارت اند از
\begin{equation}
  g_x=z_9-z_5,\qquad g_y=z_8-z_6.
  \tag{3-60}\label{eq:3-60}
\end{equation}
در نتیجه اندازه دقیق و تقریب قدرمطلق به ترتیب برابرند با
\begin{align}
  M(x,y)&=\sqrt{(z_9-z_5)^2+(z_8-z_6)^2},
  \tag{3-61}\label{eq:3-61}\\
  M(x,y)&\approx |z_9-z_5|+|z_8-z_6|.
  \tag{3-62}\label{eq:3-62}
\end{align}

عملگر سوبل برای دو مولفه گرادیان چنین نوشته می شود:
\begin{align}
  g_x&=(z_7+2z_8+z_9)-(z_1+2z_2+z_3),
  \tag{3-63}\label{eq:3-63}\\
  g_y&=(z_3+2z_6+z_9)-(z_1+2z_4+z_7).
  \tag{3-64}\label{eq:3-64}
\end{align}
ضریب ۲ به سطر یا ستون مرکزی وزن بیشتری می دهد و مقدار کمی هموارسازی درون
محاسبه مشتق وارد می کند. اندازه نهایی سوبل برابر است با
\begin{align}
M(x,y)=\Big(&[(z_7+2z_8+z_9)-(z_1+2z_2+z_3)]^2\notag\\
&+[(z_3+2z_6+z_9)-(z_1+2z_4+z_7)]^2\Big)^{1/2}.
\tag{3-65}\label{eq:3-65}
\end{align}

\BookFigure[0.84]{chapter03_figures/figure_3_51.png}
{تصویر عدسی و اندازه گرادیان سوبل. مرزها و عیب های کوچک روی محیط عدسی برجسته
شده اند.}
{fig:3-51}

\begin{latin}
\subsubsection*{Executable experiment: spatial sharpening operators}
The script compares the signed Laplacian response, Laplacian sharpening, an
unsharp mask, high-boost filtering, and Sobel magnitude on one common input.
\lstinputlisting[caption={Laplacian, unsharp, high-boost, and Sobel sharpening},label={code:sharpening}]
{chapter03_code/sharpening_filter_demo.py}
\begin{figure}[H]
  \centering
  \includegraphics[width=0.96\textwidth]{chapter03_code_outputs/sharpening_filter_result.png}
  \caption*{Executed output of the spatial sharpening script}
\end{figure}
\end{latin}

\section{ساخت بالاگذر، میان نگذر و میان گذر از پایین گذر}

شکل \ref{fig:3-52} رابطه چهار خانواده اصلی را در حوزه فرکانس نشان می دهد.
بالاگذر مکمل پایین گذر است. میان نگذر را می توان از جمع یک پایین گذر باریک و
یک بالاگذر با فرکانس قطع بزرگ تر ساخت؛ میان گذر مکمل میان نگذر است.

\BookFigure[0.84]{chapter03_figures/figure_3_52.png}
{توابع انتقال ایده آل پایین گذر، بالاگذر، میان نگذر و میان گذر. ناحیه عبور و
توقف، رفتار هسته مکانی متناظر را تعیین می کند.}
{fig:3-52}

\begin{table}[H]
\centering
\caption{ساخت چهار خانواده پالایه از هسته پایین گذر}
\label{tab:3-7}
\begin{tabular}{p{0.22\textwidth}p{0.68\textwidth}}
\toprule
نوع پالایه & هسته مکانی\\
\midrule
پایین گذر & \(lp(x,y)\)\\
بالاگذر & \(hp(x,y)=\delta(x,y)-lp(x,y)\)\\
میان نگذر & \(br(x,y)=lp_1(x,y)+[\delta(x,y)-lp_2(x,y)]\)\\
میان گذر & \(bp(x,y)=\delta(x,y)-br(x,y)\)\\
\bottomrule
\end{tabular}
\end{table}

صفحه ناحیه ای آزمایشی است که فرکانس مکانی آن با فاصله از مرکز افزایش می یابد:
\begin{equation}
  z(x,y)=\frac{1}{2}\left[1+\cos(x^2+y^2)\right].
  \tag{3-66}\label{eq:3-66}
\end{equation}
حلقه های بیرونی باریک ترند؛ بنابراین یک تصویر، مجموعه ای پیوسته از فرکانس های
کم تا زیاد را برای آزمون پالایه فراهم می کند.

\BookFigure[0.58]{chapter03_figures/figure_3_53.png}
{صفحه ناحیه ای. کاهش فاصله حلقه ها در راستای شعاع، افزایش فرکانس مکانی را
نشان می دهد.}
{fig:3-53}

\BookFigure[0.88]{chapter03_figures/figure_3_54.png}
{پروفایل یک بعدی پایین گذر و هسته دوبعدی همسانگرد حاصل از دوران آن حول مرکز.}
{fig:3-54}

\BookFigure[0.84]{chapter03_figures/figure_3_55.png}
{مقایسه پایین گذر جدایی پذیر و همسانگرد روی صفحه ناحیه ای. پاسخ چهارگوش روش
جدایی پذیر، وابستگی جهت دار آن را آشکار می کند.}
{fig:3-55}

\BookFigure[0.88]{chapter03_figures/figure_3_56.png}
{نتایج پایین گذر، بالاگذر، میان نگذر و میان گذر روی صفحه ناحیه ای، همراه با
نسخه های مقیاس شده برای نمایش.}
{fig:3-56}

\begin{latin}
\subsubsection*{Executable experiment: filter families on a zone plate}
The script generates the zone plate directly from Eq. (3-66), forms lowpass
components at two scales, and combines them to produce highpass, bandreject, and
bandpass responses.
\lstinputlisting[caption={Lowpass-derived filter families on a zone plate},label={code:zoneplate}]
{chapter03_code/zoneplate_filter_demo.py}
\begin{figure}[H]
  \centering
  \includegraphics[width=0.99\textwidth]{chapter03_code_outputs/zoneplate_filter_result.png}
  \caption*{Executed lowpass, highpass, bandreject, and bandpass results}
\end{figure}
\end{latin}

\section{ترکیب روش های بهبود مکانی}

یک روش واحد معمولا همه نیازهای تصویر را برآورده نمی کند. در نمونه پایانی،
لاپلاسین جزئیات ریز را می افزاید، گرادیان سوبل مرزهای قوی را انتخاب می کند،
هموارسازی گرادیان پیوستگی ماسک را بهتر می سازد و تبدیل توانی دامنه شدت نتیجه
را برای نمایش تنظیم می کند.

\begin{figure}[H]
  \centering
  \includegraphics[width=0.62\textwidth,height=0.72\textheight,keepaspectratio]
  {chapter03_figures/figure_3_57_part1.png}
  \caption{بخش نخست شکل ۳-۵۷: تصویر اصلی، لاپلاسین، تصویر تیزشده و گرادیان
  سوبل. هر خروجی یک نقش مشخص در زنجیره نهایی دارد.}
  \label{fig:3-57}
\end{figure}

\begin{figure}[H]\ContinuedFloat
  \centering
  \includegraphics[width=0.62\textwidth,height=0.72\textheight,keepaspectratio]
  {chapter03_figures/figure_3_57_part2.png}
  \caption{ادامه شکل ۳-۵۷: گرادیان هموارشده، ماسک حاصل ضرب، تصویر تقویت شده
  و نتیجه نهایی پس از تنظیم شدت.}
\end{figure}

\section{جمع بندی آموزشی}

این فصل یک زنجیره منطقی واحد دارد: تبدیل تک پیکسلی با \(s=T(r)\)، پردازش
آماری با هیستوگرام، پردازش همسایگی با جمع حاصل ضرب، هموارسازی با پایین گذر و
تیزسازی با مشتق. هنگام اجرای کد، دانشجو باید برای هر خروجی سه پرسش را پاسخ
دهد: کدام فرکانس ها یا شدت ها تقویت شدند، کدام جزئیات از بین رفتند، و آیا
اثر مرزی یا هاله ای ایجاد شد. پاسخ دقیق به این سه پرسش مهم تر از حفظ کردن نام
پالایه ها است.
